\chapter{Greed and Temptation}

Greed has a simple definition:

\begin{quote}
Greed is wanting something that does not belong to you.
\end{quote}

The sentence is plain, but it becomes powerful only after the word ``belong'' is understood correctly. In daily life, belonging may mean legal ownership. In investing, it also means something deeper: a return belongs to you only when your knowledge, risk control, patience, and capital structure make you worthy of receiving it.

A person who wants the result without the required ability is not merely ambitious. He may be greedy without knowing it.

This is why greed is so dangerous in capital markets. It rarely announces itself as greed. It appears as confidence, opportunity, cleverness, courage, urgency, or common sense. The person under its influence may sincerely believe that he is being practical. He may say, ``Why should I not earn more?'' He may say, ``Others are doing it.'' He may say, ``This time is different.'' He may even say, ``I deserve this.''

The market does not care what he says. It examines what he can actually carry.

\section*{The Boundary of the Word}

Before using a strong word, the boundary must be fixed. We are not trying to define greed for every possible moral, religious, political, or personal situation. Few principles work everywhere without qualification. Here the boundary is investing and wealth freedom.

Inside that boundary, greed means wanting a return, advantage, or result that is not supported by one's present competence.

This distinction matters because many people are not deliberately immoral. Their greed comes from ignorance. They want something that does not belong to them, but they do not know it does not belong to them. They imagine the result is normal. They imagine the promised return is reasonable. They imagine the risk is smaller than it is. They imagine that because the opportunity is visible, it must be available to them in the same way it is available to a trained investor.

Ignorant greed is still costly. The fact that a person did not intend to be greedy does not prevent the loss.

A simple form appears in borrowing. If one person borrows money from another and refuses to pay interest, he wants to use time, capital, and risk that belong to the lender without paying for them. Money and time have a relationship. A lender gives up present use of capital, accepts uncertainty, and waits. Interest is the price of that arrangement. Wanting the arrangement without paying the price is wanting what belongs to another person.

The same logic appears in investing. Return has a price. Sometimes the price is research. Sometimes it is volatility. Sometimes it is years of waiting. Sometimes it is the emotional discipline to hold an unpopular position. Sometimes it is the humility to refuse an opportunity one does not understand.

When the investor wants the return but refuses the price, greed has entered.

\section*{Returns That Do Not Belong to You}

A promised return is not only a number. It is also a claim about reality.

If a platform promises a very high annual return, it is claiming that somewhere behind the platform there are borrowers, businesses, trades, or projects capable of generating even more than that return after costs, defaults, fees, fraud, and operating expenses. The investor must ask whether such a claim can be true at the scale being offered.

This is where many people fail. They see only the number paid to them. They do not ask what engine must exist behind it.

Capital scale changes difficulty. Earning 10 on 100 may be easy in some small personal situation. Earning 1 billion on 10 billion is a completely different problem. The larger the capital base, the harder it becomes to obtain the same percentage return safely and repeatedly. This is why large banks usually offer low deposit rates. Their capital scale is enormous, and the money must be lent or invested through relatively controlled channels.

A compact warning table is useful:

\begin{center}
\begin{adjustbox}{max width=\linewidth}
\begin{tabular}{p{0.25\linewidth}p{0.55\linewidth}}
\hline
Return level & Question to ask \\
\hline
Low single digits & What low-risk use of capital supports it? \\
Around 15\% & What deep knowledge or durable business quality justifies it? \\
Above 25\% & What system can safely support this after costs and defaults? \\
Much higher & Am I seeing opportunity, leverage, fraud, or fantasy? \\
\hline
\end{tabular}
\end{adjustbox}
\end{center}

In the period when online peer-to-peer lending became fashionable in China, platforms that advertised annualized returns around 25\% were already displaying danger. At that level, capital rushes in. Once large amounts of capital gather, the system must find enough real projects that can use the money productively and pay more than 25\% after all costs. That is an enormous demand. In an economy growing at single-digit rates, the burden becomes even more severe.

High promised yield creates a capital-gathering effect, and gathered capital then demands real productive outlets. If those outlets do not exist in sufficient quality and scale, the system must deteriorate: defaults, Ponzi-like recycling, fraud, collapse, or some combination of these.

Many participants are not trying to steal. They simply do not know that such a return may not belong to them. They mistake availability for legitimacy. They mistake a displayed number for an earned result.

The punishment for this ignorance is usually passive suffering. The person does not know what happened until the consequence has already arrived.

\section*{The Fifteen Percent Line}

To understand what is difficult, look at people who have actually produced extraordinary long-term records.

Warren Buffett's Berkshire Hathaway produced a famous multi-decade compound annual record around 20\%. James Simons's Medallion Fund produced an even more astonishing record, though under unusual conditions and with limited outside access. Joel Greenblatt's Gotham Capital also produced spectacular returns over a long period.

These examples should not make high returns look easy. They should do the opposite. They should make the beginner humble. If the best investors in the world are admired for producing long-term compound returns in this range, then an ordinary investor should treat casual promises of 25\%, 40\%, or higher with suspicion.

Buffett once described his standard in a form worth remembering:

\begin{quote}
Buy only when there is reason to believe the stock can produce at least about 15\% compounded annually over the long term.
\end{quote}

The important number is 15\%. The more important phrase is ``compounded annually over the long term.''

Most people talk about return. Fewer talk about annualized return. Fewer still think in compound annual return. Once compounding enters the sentence, the time horizon changes. We are no longer talking about a lucky month, a good quarter, or one strong year. We are talking about many years, often ten years or more.

The formula is simple:

\[
\text{Future Value}
=
\text{Principal}
\times
(1+r)^n
\]

where \(r\) is the compound annual return and \(n\) is the number of years.

The simplicity of the formula hides the difficulty of achieving the input. A 15\% compound annual return is not a decoration. Over ten years it turns 1 into about 4.05. Over twenty years it turns 1 into about 16.37.

\begin{center}
\begin{adjustbox}{max width=\linewidth}
\begin{tabular}{rrrr}
\hline
Compound return & 5 years & 10 years & 20 years \\
\hline
10\% & 1.61x & 2.59x & 6.73x \\
15\% & 2.01x & 4.05x & 16.37x \\
25\% & 3.05x & 9.31x & 86.74x \\
\hline
\end{tabular}
\end{adjustbox}
\end{center}

The table is not a promise. It is a reality test. The higher the return, the more extraordinary the supporting engine must be. If someone offers a high return without showing the engine, the investor should not feel excited. He should feel alert.

\section*{Temptation Has Value}

Greed is not the only enemy. Temptation is subtler.

A useful definition is:

\begin{quote}
Temptation is something that has value, or appears to have value, but lacks long-term value.
\end{quote}

This definition is uncomfortable because temptation is not necessarily worthless. If it were worthless, avoiding it would be easy. Temptation often has real value. It may produce a quick gain. It may provide useful information. It may make a person feel clever. It may offer social proof. It may satisfy curiosity. It may even work once or twice.

The problem is that it pulls attention away from long-term value.

Earlier chapters argued that attention is the most important asset, that risk avoidance is the first investment necessity, and that short-term prediction is usually close to gambling. This chapter ties those ideas together. Temptation is the attractive object that consumes attention, weakens risk control, and drags the investor back into short-term prediction.

A person who decides to pursue long-term compound value will notice something strange: many things that once seemed important begin to disappear from view. Daily price movements, gossip, chart patterns, hot tips, and emotional market commentary lose some of their power. Other things become more visible: business quality, management, reinvestment ability, competitive advantage, market size, balance sheet strength, and the investor's own behavior.

The world has not changed. The filter has changed.

\begin{center}
\begin{adjustbox}{max width=\linewidth}
\begin{tikzpicture}[font=\scriptsize, >=stealth, node distance=1cm]
\node[draw, rounded corners=2pt, align=center, minimum width=2.5cm] (attention) {Attention};
\node[draw, rounded corners=2pt, align=center, above right=of attention, minimum width=2.7cm] (tempt) {Temptation\\short value};
\node[draw, rounded corners=2pt, align=center, below right=of attention, minimum width=2.7cm] (long) {Long-term\\value};
\node[draw, rounded corners=2pt, align=center, right=1.1cm of tempt, minimum width=2.5cm] (noise) {Noise and\\reaction};
\node[draw, rounded corners=2pt, align=center, right=1.1cm of long, minimum width=2.5cm] (compound) {Judgment and\\compounding};
\draw[->] (attention) -- (tempt);
\draw[->] (attention) -- (long);
\draw[->] (tempt) -- (noise);
\draw[->] (long) -- (compound);
\end{tikzpicture}
\end{adjustbox}
\end{center}

This is why the target should be stated clearly:

\begin{quote}
Pursue at least a 15\% compound annual growth target over the long term, and do so only through knowledge, discipline, and risk control.
\end{quote}

The number is not magic. It is a demanding standard. Its real function is to force the mind to ask better questions. What business could possibly support this? What price would make it reasonable? What risks could destroy the thesis? What must I know before I can deserve this return? What must I ignore so that attention remains available for what matters?

\section*{The Chart Illusion}

Almost every new investor spends time staring at price charts. This is understandable. In the beginning, there may be nothing else to look at. A chart is visible, immediate, and emotionally stimulating. It seems to show the entire story after the fact.

The beginner looks at a historical chart and thinks:

\begin{quote}
If I had sold here and bought there, I would have earned much more.
\end{quote}

This is a dangerous illusion. The chart is a record of the past, not a map handed to the past version of yourself. Looking backward, the high and low are obvious. Living forward, each price can look like both a high and a low. Every moment contains reasons to buy and reasons to sell. Every move can be explained afterward. Very few can be acted on beforehand with a reliable edge.

The chart seduces because it gives the feeling of ``what you see is what you could have done.'' That feeling is false.

Once attention is locked onto short-term charts, long-term thinking becomes difficult. The mind is pulled into the present moment again and again. Each tick demands interpretation. Each move becomes personal. The investor slowly becomes a gambler while still using investment language.

The practical danger is not the chart itself. Charts can be useful in certain methods. The danger is the beginner's relationship with the chart: stimulation mistaken for analysis, hindsight mistaken for skill, movement mistaken for meaning.

A person who cannot resist constant checking should not pretend the checking is harmless. It is training. It trains the wrong mind.

\section*{Voting Machines and Weighing Machines}

Benjamin Graham offered one of the best images in investing:

\begin{quote}
In the short run, the market is a voting machine; in the long run, it is a weighing machine.
\end{quote}

This sentence releases a great deal of unnecessary pressure.

If the market is a voting machine in the short run, then let other people vote. There is no need to vote every minute. There is no need to watch every emotional swing. There is no need to treat each temporary popularity contest as a verdict on value.

If the market is a weighing machine in the long run, then the task changes. The investor does not need to predict every vote. He needs to find what will become heavier.

In business terms, ``heavier'' means more valuable over time. More users. More revenue quality. More pricing power. More durable demand. More efficient capital allocation. More network effects. More trust. More real earnings. More ability to survive shocks and continue compounding.

The image can be made simple:

\begin{center}
\begin{adjustbox}{max width=\linewidth}
\begin{tikzpicture}[font=\scriptsize, >=stealth, node distance=1cm]
\node[draw, rounded corners=2pt, align=center, minimum width=2.7cm] (short) {Short term};
\node[draw, rounded corners=2pt, align=center, right=of short, minimum width=2.7cm] (vote) {Voting\\machine};
\node[draw, rounded corners=2pt, align=center, below=of short, minimum width=2.7cm] (long) {Long term};
\node[draw, rounded corners=2pt, align=center, right=of long, minimum width=2.7cm] (weigh) {Weighing\\machine};
\node[draw, rounded corners=2pt, align=center, right=1.1cm of vote, minimum width=2.8cm] (ignore) {Do not chase\\every vote};
\node[draw, rounded corners=2pt, align=center, right=1.1cm of weigh, minimum width=2.8cm] (study) {Study what\\gets heavier};
\draw[->] (short) -- (vote);
\draw[->] (vote) -- (ignore);
\draw[->] (long) -- (weigh);
\draw[->] (weigh) -- (study);
\end{tikzpicture}
\end{adjustbox}
\end{center}

This is the relief hidden inside long-term investing. You do not have to be good at everything. You do not have to win every short-term argument. You do not have to interpret every price movement. You have to become better at identifying durable weight and then behave in a way that lets that weight matter.

When you have learned enough to want the result with a real basis, the wanting changes. It is no longer ignorant greed. It becomes justified ambition.

\section*{How to Fight Ignorant Greed}

Because much greed comes from ignorance, the answer is not self-hatred. The answer is becoming less ignorant.

\begin{quote}
Improve a little every day. Know a little more every day.
\end{quote}

This sounds modest. It is the only reliable method.

Ignorant greed is hard to detect in the moment. If you already knew it was ignorance, it would be less dangerous. Often another person can warn you and you still cannot absorb the warning, because the missing concept is not yet alive in your mind. Only after paying a price do you see what was obvious to someone more experienced.

This is why other people's mistakes are valuable. When someone loses money chasing impossible yield, that is tuition paid in public. When someone becomes addicted to short-term charts and loses years of attention, that is instruction. When someone confuses leverage with intelligence, rumor with research, or luck with skill, he is displaying a consequence that you may be able to learn from without paying the full cost yourself.

Observation should be disciplined, not smug. The point is not to feel superior. The point is to ask:

\begin{quote}
Where am I doing the same thing in a form I have not yet recognized?
\end{quote}

A narrow checklist helps:

\begin{enumerate}
\item What exactly do I want from this decision?
\item Does this return belong to my present knowledge and risk capacity?
\item What price must be paid for this return?
\item Am I paying that price, or merely wishing to receive the result?
\item Is the value long-term, or only attractive now?
\item What evidence would show that I am being tempted?
\item What would a less ignorant version of me study before acting?
\end{enumerate}

These questions convert moral language into operating discipline.

\section*{The Real Goal}

``Getting rich'' is not a good enough goal. After food, shelter, and ordinary security are solved, money becomes complicated. For a person who has not learned how to use wealth, additional money can become a burden, a source of anxiety, a stimulant for worse desires, or a tool for magnifying old confusion.

Wealth freedom is more concrete:

\begin{quote}
From a certain point onward, you no longer have to sell your time for basic living expenses.
\end{quote}

That freedom is valuable because it gives life back to the person living it. It creates room to learn, build, help, create, think, and choose. But even wealth freedom requires a person capable of handling it. Without better concepts and better values, freedom can be wasted.

Do not spend your attention managing other people's lives in your imagination. There is a large distance between worrying about others in a vague, self-important way and actually becoming strong enough to help them. The most useful work is still your own growth: clearer concepts, stronger judgment, better habits, more accurate language, and a more honest relationship with time and capital.

Knowledge is money. Time is power.

Knowledge is money because concepts decide behavior, and behavior decides capital outcomes. Time is power because compounding needs duration, and duration rewards the person who can stay with long-term value while others chase temptation.

Greed says, ``I want the result now, whether or not I am worthy of it.''

Discipline says, ``I will become the kind of person to whom the result can belong.''

That difference is one of the central differences between a gambler and an investor, between excitement and wealth freedom, between being pulled by the market and slowly becoming free.
